\session{8}{10/31}{1限}{3}{データの収集と整備}

\goals{
  \item データの種類（構造化／非構造化、一次／二次）と主な入手先を説明できる。
  \item データ品質の観点と前処理の作業を理解し、分析できる形にデータを整えられる。
  \item 個人情報・機密情報の扱いと、AIツールへ入力する際の注意を説明できる。
}

\guidesec{解説}

\topic{データの種類}
経営で扱うデータは、いくつかの軸で分類できます。分類を知ると、「どのデータが、どの分析に、どのツールで使えるか」を判断しやすくなります。
\par\noindent\begin{gtabh}{colspec={Q[l,wd=30mm] X[l] X[l]}}
  分類の軸 & 種類 & 例 \\
  形式 & 構造化データ & 行と列で整理された表形式。売上台帳、顧客名簿、センサーの記録 \\
   & 非構造化データ & 自由記述のアンケート、メール、音声、画像、動画、契約書 \\
  収集の主体 & 一次データ & 自分（自社）が目的のために収集したデータ。販売記録、アンケート、実験 \\
   & 二次データ & 他者（政府、業界団体、企業）が収集・公開したデータ。政府統計、白書、調査報告 \\
  性質 & 定量データ & 数値で表されるデータ。売上額、人数、時間 \\
   & 定性データ & 言葉や分類で表されるデータ。顧客の声、評価、カテゴリ \\
\end{gtabh}
従来の分析は構造化データが中心でしたが、生成AIの普及で非構造化データ（文章、音声、画像）から情報を取り出すことが容易になり、
分析の対象が大きく広がりました。二次データを使うときは、収集の目的と定義（たとえば「中小企業」の定義）が自分の目的と合っているかを確認します。

\topic{データの入手先}
\begin{itemize}
  \item \textbf{社内データ}　基幹系（ERP）、顧客系（CRM）、POS、Web サイトのアクセス記録、業務日報。最も価値が高いが、部門ごとに散在していることが多い。
  \item \textbf{公開統計}　e-Stat（政府統計の総合窓口）、各府省の白書、日本銀行や業界団体の統計、RESAS（地域経済分析システム）、
        国際機関（世界銀行、OECD、国連）のデータベース。信頼性が高く、無料で使える。
  \item \textbf{Web}　企業の公式サイト、ニュース、SNS、口コミ。量は多いが、信頼性にばらつきがあり、収集（スクレイピング）には利用規約と著作権の確認が必要。
  \item \textbf{オープンデータ}　自治体や企業が再利用を認めて公開しているデータ。利用条件（ライセンス）を確認して使う。
\end{itemize}

\topic{データ品質と前処理}
集めたデータは、そのままでは分析に使えないことがほとんどです。データ品質は次の観点で評価します。
\begin{itemize}
  \item \textbf{正確性}　値が事実と一致しているか（入力ミス、計測の誤り）。
  \item \textbf{完全性}　必要な項目や期間が欠けていないか（欠損値）。
  \item \textbf{一貫性}　同じものが同じ表記・単位・定義で記録されているか（表記ゆれ、単位の混在）。
  \item \textbf{適時性}　分析の目的に対して十分に新しいか。
  \item \textbf{一意性}　同じ記録が重複していないか。
\end{itemize}
これらの問題を解消し、分析できる形に整える作業が\textbf{前処理}です。欠損値の処理（除外、補完）、重複の削除、
表記ゆれの統一（「株式会社」と「（株）」、全角と半角）、型の変換（文字列として入った日付や数値）、外れ値の確認、
1行1件・1列1項目の「整然データ」への変形などを行います。分析の作業時間の大半は前処理だと言われます。
生成AIは前処理の手順の提案とコードの生成に役立ちますが、「何を欠損とみなすか」「外れ値を除くか」は分析の目的に関わる判断であり、人が決めます。

\topic{個人情報・機密情報と AI ツール}
データを扱う際は、法令と組織の方針に従う必要があります。個人情報保護法は、個人を識別できる情報の取得・利用・提供のルールを定めています。
利用目的の特定と通知、目的外利用の禁止、第三者提供の制限、安全管理が基本です。

生成AIツールにデータを入力することは、外部の事業者のサーバーにデータを送ることを意味します。次の点を確認してください。
\begin{itemize}
  \item 入力した内容がAIの学習に使われる設定になっていないか。法人契約のツールでは学習に使わない契約が一般的だが、個人向けの無料サービスでは学習に使われることがある。
  \item 個人情報（顧客の氏名、連絡先）、機密情報（未公開の財務情報、取引条件、技術情報）は、組織の方針で許可されていない限り入力しない。
  \item 入力が必要な場合は、氏名を ID に置き換える、必要な項目だけに絞るなど、匿名化・最小化を行う。
  \item 大学の課題では、大学のAI利用規程と、この授業のAI利用ガイドに従う。他人の個人情報や企業の非公開資料は入力しない。
\end{itemize}

\topic{キーワード}
\par\noindent\begin{keywords}{}
  構造化／非構造化データ & 表形式で整理されたデータと、文章・音声・画像などそれ以外のデータ。 \\
  一次／二次データ & 自ら収集したデータと、他者が収集・公開したデータ。 \\
  e-Stat & 政府統計の総合窓口。各府省の統計を横断的に検索・取得できる。 \\
  データ品質 & 正確性・完全性・一貫性・適時性・一意性の観点で評価する。 \\
  前処理 & 欠損値の処理、重複の削除、表記ゆれの統一など、分析できる形に整える作業。 \\
  整然データ & 1行が1件、1列が1項目、1セルが1つの値になっている表の形式。 \\
\end{keywords}

\guidesec{演習課題}

\exercise{1}{AI演習：公開統計データをAIに読み込ませ、分析できる形に整える}
\instr{e-Stat から、興味のある統計表を1つダウンロードしてください（例：家計調査、小売物価統計、経済センサス）。表の先頭30行程度を次のプロンプトに貼り、構造の把握と前処理の手順を提案させてください。提案された手順のうち「自分で判断すべきこと」に印を付け、理由を書いてください。}
\begin{promptbox}[統計表の構造把握]
次の表は e-Stat からダウンロードした「（統計名）」の一部です。(1) 表の構造（見出しの行、単位、注記、集計行の有無）を説明し、(2) 1行1件・1列1項目の整然データに変換する手順を箇条書きで提案し、(3) 分析の前に確認すべきデータ品質上の懸念を挙げてください。数値を書き換えたり補ったりしないでください。
（表の先頭30行）
\end{promptbox}

\exercise{2}{データ品質のチェック}
\instr{次の顧客名簿の抜粋には、データ品質の問題が複数含まれています。問題をすべて指摘し、品質の観点（正確性・完全性・一貫性・一意性）に分類したうえで、前処理の方針を書いてください。}
\par\noindent\begin{gtab}{colspec={Q[l,wd=10mm] Q[l,wd=28mm] Q[l,wd=24mm] Q[l,wd=26mm] Q[l,wd=20mm] X[l]}}
  ID & 会社名 & 都道府県 & 登録日 & 売上（千円） & 担当 \\
  101 & 株式会社山田商事 & 愛知県 & 2025/04/01 & 1,200 & 佐藤 \\
  102 & （株）山田商事 & 愛知 & 2025-04-01 & 1200 & 佐藤 \\
  103 & 鈴木製作所 &  & 2025/13/05 & 850 & 田中 \\
  104 & 高橋物産 & 岐阜県 & 2025/05/20 & 99999 & \\
  104 & 高橋物産 & 岐阜県 & 2025/05/20 & 99999 & 田中 \\
\end{gtab}

\exercise{3}{AIツールに入力してよいか}
\instr{次の5つの場面について、大学提供の Copilot（法人契約、入力は学習に使われない設定）に入力してよいか、「可」「条件付きで可」「不可」で判定し、理由と、条件付きの場合はその条件を書いてください。}
\begin{itemize}
  \item (A) 授業で配布された公開統計の表を貼って、要約させる
  \item (B) アルバイト先の顧客名簿（氏名・電話番号つき）を貼って、重複を探させる
  \item (C) 自分が書いたレポートの下書きを貼って、誤字を指摘させる
  \item (D) 友人から見せてもらった企業のインターン選考の非公開資料を貼って、要約させる
  \item (E) アルバイト先の売上データから氏名を除き、商品別に集計した表を貼って、傾向を分析させる
\end{itemize}

\guidesec{模範解答}

\begin{answer}{演習1}
例：家計調査（二人以上の世帯、品目別支出）
\begin{itemize}
  \item 構造：先頭5行が表題・注記、6〜7行目が2段の見出し（年と月）、単位は「円」と注記にある。「総世帯」「平均」など集計行が混在。
  \item 整然データへの変換：見出しの2段を「年」「月」の列に展開し、品目を行に持つ縦長の形にする。集計行を除くか、別の列で区別する。単位の列を追加する。
  \item 品質上の懸念：「…」「-」で表された欠損・秘匿値の扱い。調査の改定（分類の変更）による期間の不連続。
\end{itemize}
自分で判断すべきこと：(1) 集計行を除くかどうか（分析の目的による）。(2) 秘匿値を欠損として除くか、ゼロとみなすか（統計の定義を確認）。
(3) 分類が変わった前後を比較してよいか。AIは手順を提案できるが、これらは統計の解説ページを読んで人が決める。
\end{answer}

\begin{answer}{演習2}
\begin{itemize}
  \item 一貫性：会社名の表記ゆれ（「株式会社」と「（株）」）、都道府県の「愛知県」と「愛知」、登録日の区切り記号（/ と -）、売上の桁区切り（1,200 と 1200）。方針：表記を統一するルールを決めて変換する。101 と 102 は同じ会社の可能性が高く、確認のうえ統合する。
  \item 完全性：103 の都道府県と、104（1行目）の担当が空欄。方針：元の記録で確認して補う。確認できなければ欠損として扱う。
  \item 正確性：103 の登録日「2025/13/05」は存在しない月。104 の売上 99999 は桁や仮入力の疑い。方針：元の伝票で確認する。確認できない場合は分析から除外し、その旨を記録する。
  \item 一意性：104 が2行あり、担当の有無だけが異なる。方針：担当が入っている行を残し、重複を削除する。
\end{itemize}
前処理の判断（統合、除外）は記録に残し、後から検証できるようにします。
\end{answer}

\begin{answer}{演習3}
\begin{itemize}
  \item (A) 可：公開データであり、個人情報・機密情報を含まない。
  \item (B) 不可：顧客の個人情報であり、アルバイト先の許可もない。重複の検索は氏名を ID に置き換えた表で行うか、職場の公式ツールで行う。
  \item (C) 可：自分の著作物で、他人の情報を含まない。ただし出力をそのまま提出しない（AI Policy）。
  \item (D) 不可：第三者の非公開資料であり、入手の経緯からも利用は不適切。
  \item (E) 条件付きで可：氏名を除いた集計表なら個人情報ではないが、売上データは職場の機密情報にあたる可能性がある。アルバイト先の責任者の許可を得ることが条件。
\end{itemize}
判断の軸は「個人情報か」「自分に利用の権限があるか」「組織の方針で許可されているか」の3つです。
\end{answer}

\guidesec{出席確認クイズの解説}
講義サイトの出席確認ページ（第8回）の5問です。
% 出席確認クイズ 第08回　データの収集と整備
%   attendance/attendance08.html から生成（解説は本ガイド用に加筆）。

\quizq{1}{「構造化データ」の例として適切なものはどれでしょうか?}%
  {音声ファイル}%
  {\correct{表形式の売上データ}}%
  {自由記述のアンケート回答}%
  {写真}%
  {正解は (b)。構造化データは行と列で整理された表形式のデータで、売上台帳や会員名簿が典型です。(a) 音声、(c) 自由記述、(d) 写真は非構造化データで、生成AIの普及によってこれらも分析しやすくなりました。}

\quizq{2}{「一次データ」と「二次データ」の説明として適切なものはどれでしょうか?}%
  {\correct{一次データは自ら収集したデータ、二次データは他者が収集・公開したデータ}}%
  {一次データは無料、二次データは有料}%
  {一次データは古く、二次データは新しい}%
  {一次データは数値、二次データは文字}%
  {正解は (a)。自らの目的のために収集したものが一次データ、政府や企業など他者が収集・公開したものを利用するのが二次データです。(b)(c)(d) の有料・新しさ・数値か文字かは、一次・二次の区別とは無関係です。二次データでは収集の目的と定義が自分の目的と合っているかの確認が必要です。}

\quizq{3}{データの前処理に含まれる作業として適切なものはどれでしょうか?}%
  {会議の設定}%
  {\correct{欠損値の処理や表記ゆれの統一}}%
  {商品の値付け}%
  {広告の作成}%
  {正解は (b)。前処理では、欠損値の処理、重複の削除、表記ゆれ（「株式会社」と「（株）」など）の統一、型の変換、外れ値の確認などを行い、分析できる状態に整えます。分析の作業時間の大半は前処理だと言われます。(a)(c)(d) はデータ整備の作業ではありません。}

\quizq{4}{日本の公的統計を横断的に入手できるポータルサイトはどれでしょうか?}%
  {\correct{e-Stat(政府統計の総合窓口)}}%
  {Googleマップ}%
  {e-Tax}%
  {マイナポータル}%
  {正解は (a)。e-Stat（政府統計の総合窓口）は、各府省が公表する政府統計を1つの窓口で検索・取得できるポータルサイトです。(c) e-Tax は国税の電子申告、(d) マイナポータルは行政手続きのサイトで、統計の入手先ではありません。}

\quizq{5}{生成AIツールに業務データを入力する際の注意として適切なものはどれでしょうか?}%
  {どんなデータでも自由に入力してよい}%
  {入力したデータは誰にも見られないので安全である}%
  {\correct{個人情報や機密情報を入力しないなど、組織の方針に従う}}%
  {個人情報を入力すると精度が上がるので推奨される}%
  {正解は (c)。生成AIツールに入力したデータは外部のサーバーに送られ、設定によっては学習に使われる可能性があるため、個人情報・機密情報の扱いは組織の方針に従います。(b)(d) は誤りで、入力したデータが「誰にも見られない」保証はなく、個人情報の入力で精度が上がるわけでもありません。}

